\documentclass[11pt]{amsart}

\usepackage{amsmath,amssymb,amsthm}
\usepackage[margin=1in]{geometry}
\usepackage[hyphens]{url}
\usepackage{hyperref}
\hypersetup{colorlinks=true,linkcolor=blue,citecolor=blue,urlcolor=blue}

\theoremstyle{plain}
\newtheorem{theorem}{Theorem}
\newtheorem{proposition}[theorem]{Proposition}
\newtheorem{lemma}[theorem]{Lemma}
\newtheorem{corollary}[theorem]{Corollary}
\theoremstyle{definition}
\newtheorem{definition}[theorem]{Definition}
\theoremstyle{remark}
\newtheorem{remark}[theorem]{Remark}

\DeclareMathOperator{\artanh}{artanh}
\newcommand{\R}{\mathbb{R}}
\newcommand{\C}{\mathbb{C}}
\newcommand{\Z}{\mathbb{Z}}
\newcommand{\N}{\mathbb{N}}

\title[Wirsching's Conjecture 3]{Wirsching's Conjecture 3, and the periodic correction to Berg--Kr\"uppel's density asymptotic}

\author{Renato Augusto Tavares}
\address{Universidade Federal de Goi\'as}
\email{dr.renatotavares@gmail.com}

\subjclass[2020]{11B83, 39B22, 30D05, 11M06}
\keywords{Collatz map, predecessor density, invariant density, log-periodic asymptotics,
Mellin transform, atomic function}

\begin{document}

\begin{abstract}
Wirsching's 2003 argument for positive predecessor density under the $3n+1$ map reduces to three
conjectures. Conjecture 3 concerns the near-zero asymptotics of an invariant density $\varphi$,
compared against an explicit asymptotic $\varphi_0$ due to Berg and Kr\"uppel. The exact
log-Laplace transform of $\varphi$ decomposes as a smooth part plus a periodic correction $H$. Berg
and Kr\"uppel gave this correction, for this eigenfunction, as an infinite product in 1998, five
years before Wirsching's conjecture; they did not ask whether it is constant. We give $H$ instead
as a Fourier series with closed-form coefficients in $\Gamma$ and $\zeta$, and certify rigorously
that it is not. Wirsching's own comparison class fixes a single phase of $H$; there, Conjecture 3
holds, with an explicit limit. Away from that class the phase sweeps a full period, so the
stronger, unrestricted asymptotic $\varphi(t) \sim c\,\varphi_0(t)$ as $t \to 0^+$ fails. Wirsching
needs less than the literal statement of Conjecture 3; we check directly that his actual
requirement holds with a wide margin.
\end{abstract}

\maketitle

\section{Introduction}

Let $T$ denote the $3n+1$ map on the positive integers, $T(n) = n/2$ for $n$ even and $T(n) =
(3n+1)/2$ for $n$ odd. Wirsching's monograph \cite{wirsching1998} develops the predecessor-set
machinery this paper's argument descends from; \cite{wirsching2003} takes it further, studying the
density of predecessor sets under $T$ through an averaging construction on the $3$-adic integers: a
family of operators $S_l$, built from path counts in the predecessor graph, converges in the strong
operator topology to a limiting operator $S_\infty$ whose essential part is a Markov chain on
$[0,1]$. That chain has a unique invariant density $\varphi$, and Wirsching's proof of positive
predecessor density reduces to three conjectures about $\varphi$ and related quantities, stated as
three unbridged gaps in the argument.

The load-bearing one is Conjecture 3. It concerns the behaviour of $\varphi$ as its argument tends
to $0$, compared against an explicit comparison function $\varphi_0$. Berg and Kr\"uppel
\cite{bergkruppel1998} had already produced $\varphi_0$ five years earlier, by Laplace transform
and saddlepoint analysis, as the asymptotic solution of a truncated version of the functional
equation $\varphi$ satisfies. Their argument for the asymptotic comparability of $\varphi$ and
$\varphi_0$ was informal, and left open both the exact multiplicative constant and whether the
ratio converges at all. Wirsching considers the strongest possible reading of that gap directly:

\begin{quote}
``That replacement would be possible, if, e.g., $\varphi \sim c\,\varphi_0$ for some constant
$c > 0$. It turns out that the following (slightly weaker) conjecture \dots suffices.''
\end{quote}

He then states Conjecture 3: along a specific comparison class of sequences $(z_l)$ tending to $0$
at the borderline rate $z_l \sim l \cdot 3^{-l}$, the ratio $\varphi(z_l)/\varphi_0(z_l)$ converges
to a positive constant, uniformly on the class. This is weaker than the unrestricted asymptotic he
first considers. The constant cancels out of the one inequality his own proof actually needs, so a
class-dependent limit is enough.

Neither statement had been settled. This paper settles both, in opposite directions.

\begin{theorem}[Conjecture 3]
\label{thm:conj3}
Wirsching's Conjecture 3 is true. On his comparison class $\widetilde{A}_\delta$, the ratio
$\varphi(z_l)/\varphi_0(z_l)$ converges uniformly to $e^{H(0)} = 0.534122\ldots$, where $H$ is the
periodic function of Theorem~\ref{thm:H} below and the normalization of $\varphi_0$ is Berg and
Kr\"uppel's own, equation~(9.6) of \cite{bergkruppel1998}.
\end{theorem}

\begin{theorem}[Failure of the unrestricted asymptotic]
\label{thm:unrestricted}
The stronger asymptotic $\varphi(t) \sim \kappa\,\varphi_0(t)$ as $t \to 0^+$, for any single
constant $\kappa$ and with no restriction to a comparison class, is false. The ratio
$\varphi(t)/\varphi_0(t)$ oscillates indefinitely as $t \to 0^+$, with period $\log 3$ in
$\log(1/t)$ and amplitude bounded below by a positive, certified constant.
\end{theorem}

The two theorems have a single common source. Write $t = e^{-\tau}$ and let $H(w)$ be the
$\log 3$-periodic function appearing in the exact decomposition of $\log\varphi$'s log-Laplace
transform (Theorem~\ref{thm:H}). Wirsching's comparison class fixes $\tau$'s phase modulo $\log 3$
at a single value; every other class of comparison sequences sweeps through all phases. Whether the
ratio converges or oscillates depends on which sequences one compares along, and the two theorems
above read the same object at two different scopes.

\begin{corollary}
\label{cor:bk}
$H$ is Berg and Kr\"uppel's own periodic factor for the eigenfunction $\varphi$, from
Proposition~9.3 of \cite{bergkruppel1998}, which they give as an infinite product. Its Fourier
series, Proposition~\ref{prop:Hfourier} below, is new.
\end{corollary}

Section~\ref{sec:prelim} recalls the functional equation $\varphi$ satisfies and Berg and
Kr\"uppel's comparison asymptotic. Section~\ref{sec:periodic} constructs $H$, proves
Proposition~\ref{prop:Hfourier} and Corollary~\ref{cor:bk}, and certifies that $H$ is not constant.
Sections~\ref{sec:formulaA} and~\ref{sec:envelope} transfer this from the log-Laplace transform to
$\varphi$ itself, through a uniform saddlepoint approximation and an envelope estimate.
Section~\ref{sec:bk} identifies the resulting smooth part with Berg and Kr\"uppel's saddlepoint
expression and pins down every constant. Section~\ref{sec:proofs} proves both theorems and checks
Wirsching's actual requirement directly. Section~\ref{sec:discussion} states what remains open.

\section{The invariant density and its Laplace transform}
\label{sec:prelim}

Let $U_1, U_2, \dots$ be independent, identically distributed uniform random variables on $[0,1]$,
and set
\[
X = \sum_{j \ge 1} \frac{2}{3^j} U_j \in [0,1].
\]
The self-similar identity $X = \tfrac{2}{3}U_1 + \tfrac{1}{3}X'$, with $X'$ an independent copy of
$X$, shows that the density $\varphi$ of $X$ satisfies
\[
\varphi(x) = \frac{3}{2}\int_{3x-2}^{3x} \varphi,
\]
equivalently $\varphi'(x) = \tfrac{9}{2}\varphi(3x)$ on $(0, \tfrac{2}{3})$: this is Wirsching's
invariant density (up to the reflection $x \mapsto 1-x$ fixing the convention), a translate of
Rvachev's atomic function $h_3$ \cite{rvachev1990}, and the dilation-$3$, $\lambda=2/3$ member of
Berg and Kr\"uppel's family of functional equations $\lambda\, g'(t) = 3\bigl(g(3t) -
g(3t-2)\bigr)$. It is $C^\infty$, supported on $[0,1]$, and constant, equal to $3/2$, on the middle
third $[1/3, 2/3]$.

For $s > 0$ write
\[
K(s) = \log \mathbb{E}\bigl[e^{-sX}\bigr] = \sum_{j \ge 1} g\!\left(\frac{2s}{3^j}\right), \qquad
g(b) = \log\frac{1-e^{-b}}{b}.
\]
The branch of $g$ is the one analytic on $\{\operatorname{Re} w > 0\}$ and real on $(0,\infty)$;
since $(1-e^{-w})/w$ is analytic and non-vanishing there (its zeros lie on $2\pi i \Z \setminus
\{0\}$, and $w=0$ is a removable singularity with value $1$), and the half-plane is simply
connected, $g$ exists and is unique. Differentiating termwise, justified because $g(w) = -w/2 +
O(w^2)$ as $w \to 0$ so the tail of the series is normally convergent together with every
derivative,
\[
t(s) := -K'(s), \qquad V(s) := s^2 K''(s) > 0,
\]
the second an exact positivity, $K''(s)$ being the variance of $X$ under its exponential tilt by
$e^{-sX}$. As $s$ ranges over $(0,\infty)$, $t(s)$ decreases strictly from $1/2$ to $0$ (proved in
Lemma~\ref{lem:saddles} below), so studying $\varphi(t(s))$ as $s \to \infty$ is the same as
studying $\varphi(t)$ as $t \to 0^+$.

Berg and Kr\"uppel's comparison function solves the truncated equation $\lambda\, g'(t) = d\, g(dt)$
by Laplace transform and saddlepoint analysis. In their normalization (their equation~(9.6), with
$\alpha = \tfrac12 - \tfrac{\log 2}{\log 3}$, $\beta = \tfrac{1}{2\log 3}$, and $\gamma$,
$\delta_{\mathrm{BK}}$, $\varepsilon$ determined by $\alpha$ and $\beta$ through their
Proposition~9.1),
\[
\varphi_0(t) = \frac{(2\beta)^\varepsilon}{\sqrt{2\pi}}\; t^\gamma\,(-\log t)^{\delta_{\mathrm{BK}}}
\, \exp\!\left(-\beta \log^2\!\frac{t}{-\log t}\right).
\]
Write $\varphi_{0,\mathrm{bare}}$ for the same expression without the leading constant
$(2\beta)^\varepsilon/\sqrt{2\pi}$: the two differ only by that constant.

\section{The periodic correction}
\label{sec:periodic}

Set $c = \log 3$ and write $w = \log s$, $L(w) = K(e^w)$. Substituting into the defining series and
splitting off the $j=1$ term against the rest reproduces the recurrence
\begin{equation}
\label{eq:recurrence}
L(w+c) - L(w) = \log\bigl(1 - e^{-2e^w}\bigr) - w - \log 2.
\end{equation}
Let
\[
Q(w) = -\frac{w^2}{2c} + \left(\frac12 - \frac{\log 2}{c}\right) w,
\]
chosen so that $Q(w+c) - Q(w) = -w-\log2$, matching \eqref{eq:recurrence}'s non-periodic terms
exactly. Write $a := \tfrac12 - \tfrac{\log 2}{c}$ for $Q$'s linear coefficient, used again from
Section~\ref{sec:envelope} onward.

\begin{theorem}
\label{thm:H}
Define
\begin{equation}
\label{eq:Hexact}
H(w) := L(w) - Q(w) + \sum_{k \ge 0} \log\bigl(1 - e^{-2\cdot 3^k e^w}\bigr).
\end{equation}
Then $H$ is exactly $c$-periodic, $H(w+c) = H(w)$ for every $w \in \R$, and
\[
L(w) = Q(w) + H(w) + \Delta(w), \qquad \Delta(w) = -\sum_{k\ge 0}\log\bigl(1-e^{-2\cdot3^k e^w}\bigr),
\]
with $\Delta^{(j)}(w) = O_j\bigl(e^{jw - 2e^w}\bigr)$ for every $j \ge 0$: doubly exponentially
small, together with all its derivatives, as $w \to +\infty$.
\end{theorem}

\begin{proof}
Write $d(w) = \log(1-e^{-2e^w})$, so \eqref{eq:recurrence} reads $L(w+c)-L(w) = d(w) - w - \log 2$.
Directly from the definition of $Q$ (expand $(w+c)^2 = w^2+2wc+c^2$ and collect terms),
$Q(w+c)-Q(w) = -w-\log2$ as well, so the two non-periodic parts cancel exactly:
\[
(L-Q)(w+c) - (L-Q)(w) = \bigl[d(w)-w-\log2\bigr] - \bigl[-w-\log2\bigr] = d(w).
\]
The series in \eqref{eq:Hexact} telescopes this exactly: with $D(w) := (L-Q)(w)$,
\[
D(w) + \sum_{k\ge0} d(w+kc) = D(w+c) - d(w) + \sum_{k\ge0}d(w+kc) = D(w+c) + \sum_{k\ge0}d(w+(k+1)c),
\]
so the right side of \eqref{eq:Hexact} is invariant under $w \mapsto w+c$, that is, $H$ is
$c$-periodic. Convergence of the series and its termwise differentiability, together with the
stated bound on $\Delta^{(j)}$, follow from $d(w) = O(e^{-2e^w})$ and the corresponding bounds on
its derivatives, since $2\cdot3^k e^w \ge 2 e^w$ for $k \ge 0$ makes the series dominated, uniformly
for $w$ in any half-line bounded away from $-\infty$, by a geometric series in $3^{-k}$ inside the
exponent.
\end{proof}

The decomposition $L = Q+H+\Delta$, with $\Delta$ negligible as $w\to\infty$, is unique: any two
periodic functions differing from $L-Q$ by an $o(1)$ term as $w\to\infty$ must be equal, being
periodic functions that agree in the limit along every residue class of $w$ modulo $c$.

\begin{remark}
Berg and Kr\"uppel's own Proposition~9.3 (specialized to $a=3$, $b=3/2$ in their notation, exactly
this eigenfunction) gives the equivalent identity $\exp(H(w)) = a^{\frac12 t^2 - \alpha t}
\prod_{k\ge0}\frac{1-e^{-a^{t-k}/b}}{a^{t-k}/b}\prod_{l\ge1}\bigl(1-e^{-a^{t+l}/b}\bigr)$, $t=w/c$,
as an infinite product, in their proof of that proposition. The Fourier series of
Proposition~\ref{prop:Hfourier} below is what lets us certify non-constancy (Proposition
\ref{prop:Hcert}), which their product form does not make transparent.
\end{remark}

\begin{proposition}
\label{prop:Hfourier}
$H$ has the Fourier expansion $H(w) = \sum_{m\in\Z} \hat H(m) e^{i\omega_m w}$, $\omega_m =
2\pi m/c$, with
\[
\hat H(0) = \frac{\log 2}{2} - \frac{\log^2 2}{2c} - \frac{c}{12} - \frac{A}{c}, \qquad
A = \frac{\pi^2}{12} - \frac{\gamma_E^2}{2} - \gamma_1,
\]
\[
\hat H(m) = -\frac{2^{i\omega_m}}{c}\, \Gamma(-i\omega_m)\, \zeta(1-i\omega_m), \qquad m \ne 0,
\]
where $\gamma_E$ is the Euler--Mascheroni constant and $\gamma_1$ the first Stieltjes constant.
\end{proposition}

\begin{proof}[Sketch\footnote{Proofs marked ``Sketch'' give the complete argument at the level of
a referee report; every numerical constant they produce is independently certified in the
accompanying repository (Section~\ref{sec:code}), and full line-by-line derivations are on record
in the underlying project's notes, cited there.}]
The Mellin transform of $g$ on the strip $-1 < \operatorname{Re}\, z < 0$ is $g^*(z) =
-\Gamma(z)\zeta(1+z)$, obtained by integrating by parts onto $g'(b) = 1/(e^b-1) - 1/b$ and using the
classical continuation $\int_0^\infty\bigl(\tfrac{1}{e^u-1}-\tfrac1u\bigr)u^{z-1}\,du =
\Gamma(z)\zeta(z)$. Since $K(s) = \sum_j g(2s/3^j)$, the Mellin transform of $s\mapsto K(s)$ picks
up a factor $2^z\sum_j 3^{-jz} = 2^z/(3^z-1)$ from this harmonic sum; shifting the resulting
inversion contour past the triple pole at the origin and past each simple pole at $z = i\omega_m$
(where $3^z=1$) reproduces $Q + \hat H(0)$ and the Fourier modes above, respectively, the factor
$2^{i\omega_m}$ coming from that same harmonic-sum step.
\end{proof}

\begin{remark}
This is an instance of the de Bruijn--Mahler phenomenon: a base-change harmonic sum, Mellin
transformed, picks up simple poles on a vertical line, and their residues assemble into an exactly
periodic correction with Fourier coefficients in $\Gamma$ and $\zeta$. De Bruijn found the first
case of this in Mahler's partition asymptotic \cite{debruijn1948}; Erd\H{o}s and Richmond named it
\cite{erdosrichmond1976}; Kato and McLeod \cite{katomcleod1971} and Derfel \cite{derfel1995}
established the log-squared decay rate of \eqref{eq:Hexact}'s companion asymptotic $\varphi_0$ as
generic across this whole class of rescaled functional equations. Our contribution here is the
closed form for the periodic correction itself, not the mechanism producing it.
\end{remark}

The series for $H'$ and $H''$ converge absolutely and uniformly, since $|\Gamma(-i\omega_m)| =
\sqrt{\pi/(\omega_m \sinh \pi\omega_m)}$ decays exponentially in $m$ while $\zeta(1-i\omega_m)$
grows only polynomially, and this yields explicit, certified bounds:
\begin{equation}
\label{eq:Hbounds}
\sup_w |H'(w)| \le 0.0011977472315550332, \qquad \sup_w |H''(w)| \le 0.0068518962896650951.
\end{equation}

\begin{proposition}[Certified non-constancy]
\label{prop:Hcert}
$H$ is not constant. Rigorously,
\[
-0.000377190280943987 \;<\; H(0) - H(\log(3/2)) \;<\; -0.000377190280943985,
\]
and consequently $\operatorname{osc}(H) := \sup H - \inf H \ge 3.7719\times 10^{-4}$; a full
optimization gives the certified enclosure
$\operatorname{osc}(H) \in [4.18744947692\times10^{-4},\, 4.18756644224\times10^{-4}]$.
\end{proposition}

\begin{proof}[Sketch]
Write $g(b) = -b + \log S(b)$, $S(b) = \sinh(b)/b = \sum_{n\ge0} b^{2n}/(2n+1)!$. Termwise
domination of $6^n n! \le (2n+1)!$ gives $0 \le \log S(b) \le b^2/6$ for $b \ge 0$, hence an
explicit, two-sided enclosure of the tail $\sum_{r>R} g(s/3^r)$ for any $R$ and $s > 0$; combined
with a finite head evaluated in interval or ball arithmetic (via \eqref{eq:Hexact} directly, or
independently via ball arithmetic in $\Gamma$ and $\zeta$ applied to Proposition
\ref{prop:Hfourier}), this certifies the stated interval to arbitrary precision. The full
oscillation is certified the same way, via the Lipschitz bound in \eqref{eq:Hbounds} and a fine
grid search over one period.
\end{proof}

\section{A uniform saddlepoint approximation}
\label{sec:formulaA}

Define, for $s>0$,
\[
\varphi_{\mathrm{sp}}(t(s)) := \frac{\exp\bigl(K(s)+st(s)\bigr)}{\sqrt{2\pi K''(s)}}.
\]
This section proves that $\varphi(t(s))/\varphi_{\mathrm{sp}}(t(s)) \to 1$ as $s \to \infty$,
uniformly over every phase of $s$ modulo $3^{\Z}$, which is what makes the two theorems of
Section~\ref{sec:proofs} depend on different scopes rather than on different mathematics.

For $n \ge 2$, define the derivative-normalized functions $\kappa_n(w) = (-1)^n w^n g^{(n)}(w)$ on
$\{\operatorname{Re} w > 0\}$; $\kappa_n(b)$, $b>0$, is $b^n$ times the $n$-th cumulant of the
random variable with density $b\,e^{-bv}/(1-e^{-b})$ on $[0,1]$. Writing $x=w/2$,
\[
\kappa_2(w) = 1 - \frac{x^2}{\sinh^2 x}, \qquad
\kappa_3(w) = 2 - \frac{2x^3\cosh x}{\sinh^3 x}, \qquad
\kappa_4(w) = 6 + \frac{2x^4(1-3\coth^2 x)}{\sinh^2 x},
\]
and from the Bernoulli expansion $g(w)+w/2 = \sum_{k\ge1} B_{2k}w^{2k}/(2k(2k)!)$ on $|w|<2\pi$,
\[
\kappa_n(w) = (-1)^n\!\!\sum_{k\ge\lceil n/2\rceil}\!\! \frac{B_{2k}\,w^{2k}}{2k\,(2k-n)!}, \qquad n\ge2,
\]
so that $\kappa_n(w) = O(w^2)$ as $w \to 0$ for every $n \ge 2$.

Write $s = \rho\cdot 3^N/2$ for the unique $N = \lfloor \log_3(2s)\rfloor \in \Z$ and $\rho \in
[1,3)$, and $b_j = 2s/3^j$. Then $\{b_j : j \ge 1\} = \{\rho\cdot 3^m : m \le N-1\}$: exactly $N$ of
these, the smallest being $\rho \ge 1$, are $\ge 1$; the rest form a geometric tail below $1$. This
single structural fact is the source of every phase-independent constant below.

\begin{lemma}[Boundedness off the real axis]
\label{lem:kappa}
For $w = b(1+iu)$, $b>0$, $|u|\le 1$, with $e(r) = r^2/\sinh^2 r$ and $f(r) = r^3\cosh r/\sinh^3 r$,
\[
|\kappa_2(w)| \le 1 + 2e(b/2), \qquad |\kappa_3(w)| \le 2 + 2\cdot 2^{3/2} f(b/2).
\]
Both bounds are finite over the whole sector $|\operatorname{Im}w|\le\operatorname{Re}w$; $M_2 :=
\sup |\kappa_2| \le 3$, $M_3 := \sup|\kappa_3| \le 2+2^{5/2} = 7.657$. In addition, for $|w|\le2$,
$|\kappa_2(w)| \le 0.114|w|^2$ and $|\kappa_3(w)|\le 0.0119|w|^4$.
\end{lemma}

\begin{proof}
$|\sinh z|^2 = \sinh^2(\operatorname{Re}z)+\sin^2(\operatorname{Im}z) \ge \sinh^2(\operatorname{Re}z)$
and $|\cosh z|^2 \le \cosh^2(\operatorname{Re}z)$; the first excludes the poles of $\coth$, since
$\operatorname{Re}(w/2)=b/2>0$ never approaches $\pi i k$. This gives the sector bounds directly.
Monotonicity of $e$ is immediate from the increasing power series of $\sinh(r)/r$; monotonicity of
$f$ is Lazarevi\'c's inequality, $(\sinh r/r)^3 > \cosh r$, checked by comparing alternating
truncated series near $r=0$ and by monotonicity of $\coth r - 1/r$ away from it. The local bounds
at $|w|\le2$ follow by applying the maximum-modulus principle to $\kappa_2(w)/w^2$ and
$\kappa_3(w)/w^4$ (analytic on $|w|<2\pi$ by the Bernoulli expansion): bound each on $|w|=2$ by the
triangle inequality, using $|B_{2k}| = 2(2k)!\,\zeta(2k)/(2\pi)^{2k}$ and $\zeta(2k)\le\zeta(2)$,
respectively $\zeta(4)$.
\end{proof}

\begin{lemma}[Variance]
\label{lem:variance}
For all $\rho\in[1,3)$, $N\ge1$, with $V := s^2K''(s) = \sum_j\kappa_2(b_j)$,
\[
N - A_V \le V \le N + 0.1283, \qquad A_V := \sum_{m\ge0} e(3^m/2) = 1.4269413069\ldots.
\]
\end{lemma}

\begin{proof}
Split $V = \sum_{m=0}^{N-1}\kappa_2(\rho3^m) + \sum_{m<0}\kappa_2(\rho3^m)$, using $\kappa_2(b) =
1-e(b/2) \in (0,1)$. For the lower bound, drop the second (non-negative) sum and use $e(\rho3^m/2)
\le e(3^m/2)$ ($e$ decreasing, $\rho\ge1$) in the first, giving $\sum_{m=0}^{N-1}\kappa_2(\rho3^m)
\ge N - A_V$. For the upper bound, the first sum is $<N$ term by term, and the second sum, where
$\rho3^m<1$, is bounded by Lemma~\ref{lem:kappa}'s local estimate: $\sum_{m<0}\kappa_2(\rho3^m) \le
0.114\rho^2\sum_{m<0}9^m = 0.114\rho^2/8 \le 0.1283$ for $\rho<3$.
\end{proof}

\begin{lemma}[Tail bound and existence of the density]
\label{lem:tail}
If $b_0\ge1$ and $b_k=3^kb_0$, then $\sum_{k\ge0}\log\coth(b_k/2) \le 3e^{-b_0}$. Consequently, with
$\Lambda(y) := K(s(1+iy)) - K(s) + isy\,t(s)$,
\[
|e^{\Lambda(y)}| \le e^{3/e}(1+y^2)^{-N/2} \le 3.0152\,(1+y^2)^{-N/2}
\]
for every real $y$. In particular $X$ has a continuous density, the tilted characteristic function
is integrable for $N\ge3$, and Fourier inversion at $x=t(s)$ gives
\begin{equation}
\label{eq:R}
R(s) := \frac{\varphi(t(s))}{\varphi_{\mathrm{sp}}(t(s))} = \sqrt{\frac{V}{2\pi}}\int_\R e^{\Lambda(y)}\,dy,
\end{equation}
a real quantity, since $\Lambda(-y) = \overline{\Lambda(y)}$.
\end{lemma}

\begin{proof}
$\log\coth(x/2) = 2\artanh(e^{-x}) \le 2e^{-x}/(1-e^{-2x})$; summing over the geometric orbit
$b_k=3^kb_0$ and bounding the resulting series by its first few terms gives the stated constant.
$X$'s density is continuous because the sum of the first two summands already has a continuous,
compactly supported density and convolution preserves continuity. For $\Lambda$: $t(s) =
-K'(s)$ cancels the linear term of the tilted transform exactly, and each remaining factor
$e^{g(b_j(1+iy))-g(b_j)}$ has modulus at most $1$ (a tilted characteristic function), with modulus
at most $\coth(b_j/2)/\sqrt{1+y^2}$ on the $N$ indices where $b_j\ge1$; the tail estimate applies to
exactly that geometric set, with $b_0=\rho\ge1$.
\end{proof}

\begin{lemma}[Central expansion]
\label{lem:taylor}
With $h(u) = g(b(1+iu))$, Taylor's theorem in \emph{integral} form (the Lagrange form is invalid
for a complex-valued function of a real variable) gives
\[
h(y) = h(0)+h'(0)y+h''(0)\tfrac{y^2}{2} + \tfrac12\int_0^y (y-u)^2 h'''(u)\,du.
\]
Consequently, for $\theta\in(0,1]$,
\[
S := \sum_j \sup_{|u|\le \theta}|h_j'''(u)| \le 2N + 3.7442,
\]
and hence $|\Lambda(y) + Vy^2/2| \le B|y|^3$ for $|y|\le \theta$, with $B := S/6 \le (2N+3.7442)/6$.
\end{lemma}

\begin{proof}
Write $h_j(u) := g(b_j(1+iu))$ for each $j$, so $h_j'''(u) = \kappa_3(b_j(1+iu))/(1+iu)^3$ up to a
sign. The bound on $h_j'''$ follows from Lemma~\ref{lem:kappa} applied to $\kappa_3$, summed over
the $N$ indices with $b_j\ge1$ (via $f$'s monotonicity, worst case $\rho=1$) and the remaining
geometric tail (via the local bound).
\end{proof}

\begin{theorem}[Formula (A)]
\label{thm:formulaA}
For $N\ge17$, uniformly over $\rho\in[1,3)$, $|R(s)-1|\le E(N) := e_1(N)+e_2(N)+e_3(N)$, with
$\theta_N := N^{-1/3}$, $V_{\mathrm{lo}} := N-A_V$, $V_{\mathrm{up}} := N+0.1283$,
$B := (2N+3.7442)/6$,
\begin{align}
e_1(N) &:= \frac{4\,e^{B\theta_N^3}\,B}{\sqrt{2\pi}\,V_{\mathrm{up}}^{3/2}}, \qquad
e_2(N) := \frac{2}{\sqrt{2\pi}}\cdot\frac{e^{-\theta_N^2V_{\mathrm{lo}}/2}}{\theta_N\sqrt{V_{\mathrm{lo}}}},
\label{eq:Ebound}\\
e_3(N) &:= 2\,e^{3/e}\sqrt{\frac{V_{\mathrm{up}}}{2\pi}}\cdot
\frac{(1+\theta_N^2)^{-(N-2)/2}}{\theta_N(N-2)}. \notag
\end{align}
$E$ is strictly decreasing, $E(17)<1$, $E(N)\le4.1\,N^{-1/2}$ for all $N\ge17$, and $\sqrt N\,E(N)
\to \tfrac43e^{1/3}(2\pi)^{-1/2} = 0.742358\ldots$. Consequently $\varphi(t) =
\varphi_{\mathrm{sp}}(t)\,(1+o(1))$ as $t \to 0^+$, with no restriction on the sequence of $t$'s.
\end{theorem}

\begin{proof}
Split \eqref{eq:R} at $|y|=\theta_N$. On $|y|\le\theta_N$, Lemma~\ref{lem:taylor} bounds the cubic
remainder $|\Lambda(y)+Vy^2/2|$ by $B\theta_N^3$, which stays bounded (not $o(1)$: this is exactly
why $e_1$ carries the factor $e^{B\theta_N^3}$ rather than vanishing) as $N\to\infty$, since $B=O(N)$
and $\theta_N^3=N^{-1}$. Writing $C(y):=e^{\Lambda(y)}-e^{-Vy^2/2}$, $|e^z-1|\le|z|e^{|z|}$ gives
$|C(y)|\le B|y|^3e^{B\theta_N^3}$ on this range, and $\int_\R e^{-Vy^2/2}|y|^3\,dy=4/V^2$ (using
$V\ge V_{\mathrm{lo}}$, from Lemma~\ref{lem:variance}) gives the central Taylor-remainder
contribution $e_1$. The truncated Gaussian tail $2\int_{\theta_N}^\infty
\tfrac1{\sqrt{2\pi}}e^{-Vy^2/2}\,dy$, bounded via $\int_c^\infty e^{-v^2/2}\,dv\le e^{-c^2/2}/c$ at
$c=\theta_N\sqrt{V_{\mathrm{lo}}}$, gives $e_2$. On $|y|>\theta_N$, Lemma~\ref{lem:tail}'s bound
$|e^{\Lambda(y)}|\le e^{3/e}(1+y^2)^{-N/2}$ integrates, via $\int_a^\infty(1+y^2)^{-N/2}\,dy \le
(1+a^2)^{-(N-2)/2}/(a(N-2))$ at $a=\theta_N$ (using $V\le V_{\mathrm{up}}$ for the leading
$\sqrt{V/2\pi}$ factor in \eqref{eq:R}), to give $e_3$; this needs $\theta_N \gg N^{-1/2}$, satisfied
since $\tfrac13<\tfrac12$. Every constant entering $E$ (via $A_V$, Lemma~\ref{lem:variance}, and the
$3.0152\approx e^{3/e}$ of Lemma~\ref{lem:tail}) is a supremum over $b>0$ or a sum over one complete
geometric orbit bounded using $\rho\ge1$, hence independent of $\rho$. Evaluating \eqref{eq:Ebound}
numerically gives the stated $N_0=17$, monotonicity, and limit. Since $t(s)$ ranges bijectively over
$(0,1/2)$, this is exactly the statement $\varphi(t)=\varphi_{\mathrm{sp}}(t)(1+o(1))$ as $t\to0^+$.
\end{proof}

\begin{remark}
\label{rem:edgeworth}
The rate is not sharp. Carrying the central expansion to fourth order, using that the odd cubic
term integrates to zero against the Gaussian core, and using the limits $V_3 := \sum_j\kappa_3(b_j)
= 2N+O(1)$, $V_4 := \sum_j\kappa_4(b_j) = 6N+O(1)$ (the third and fourth cumulant sums, by the same
argument as Lemma~\ref{lem:variance}) gives the Edgeworth correction $R(s) = 1 - \tfrac1{12N} +
O(N^{-3/2})$,
matching the classical Stirling correction for a $\mathrm{Gamma}(N,1)$ law: after tilting, the
$b_j\ge1$ summands behave asymptotically like independent $\mathrm{Exp}(1)$ variables, so $X$ under
its tilt is asymptotically $\mathrm{Gamma}(N,1)$-distributed, and $\rho$ enters only through
$O(1)$ shifts of the cumulants, at order $N^{-2}$. We do not need this sharper rate and do not
prove the remainder term here; Theorem~\ref{thm:formulaA}'s qualitative $o(1)$ is all that
Section~\ref{sec:proofs} uses.
\end{remark}

\section{The envelope lemma}
\label{sec:envelope}

Theorem~\ref{thm:formulaA} compares $\varphi$ to $\varphi_{\mathrm{sp}}$, a quantity built from the
full transform $K=Q+H+\Delta$. It remains to compare $\varphi_{\mathrm{sp}}$ to a smooth auxiliary
system built from $Q$ alone, so that the periodic part $H$ can be extracted explicitly. Write
$t=e^{-\tau}$, $g_\tau(w) = L(w)+e^{w-\tau}$, $g_{0,\tau}(w) = Q(w)+e^{w-\tau}$. The true and smooth
systems each have their own analogue of $-K'$: write $B_{\mathrm{tr}}(w) := e^wt(e^w) = -L'(w)$ and
$B_{\mathrm{sm}}(w) := -Q'(w) = w/c-a$.

\begin{lemma}[True and smooth saddles]
\label{lem:saddles}
For each $\tau$, $g_\tau$ has a unique global minimizer $w^*(\tau)$, at which $g_\tau''(w^*) =
V(r^*)>0$, $r^*=e^{w^*}$; and $g_{0,\tau}$ has a unique local minimizer $w_0(\tau)$ among solutions
with $B_0 := B_{\mathrm{sm}}(w_0) = e^{w_0-\tau} > 1/c$. Both are characterized by $\tau =
\Phi(w^*) = \Phi_0(w_0)$ for $\Phi(w) = w-\log B_{\mathrm{tr}}(w)$ (a strictly increasing bijection
$\R\to(\log2,\infty)$, since $t(s)$ is strictly decreasing) and $\Phi_0(w)=w-\log B_{\mathrm{sm}}(w)$,
and satisfy
\[
|w^*-w_0| = O(1/B_0), \qquad 0 \le g_\tau(w_0)-g_\tau(w^*) = O(1/B_0).
\]
\end{lemma}

\begin{proof}[Sketch]
Strict convexity of $K$ gives the global uniqueness of $w^*$. At $w_0$, $g_\tau'(w_0) = H'(w_0) +
O(e^{-2e^{w_0}})$ by Theorem~\ref{thm:H}, bounded by \eqref{eq:Hbounds}; since $g_\tau''$ is bounded
above and below by explicit constants on $[w_0-1,w_0+1]$ (again via \eqref{eq:Hbounds} and
Lemma~\ref{lem:kappa}), a sign-change argument for $g_\tau'$ locates $w^*$ within $O(1/B_0)$ of
$w_0$, and Taylor's theorem about $w^*$ then bounds the resulting gap in $g_\tau$-values by the
same order.
\end{proof}

\begin{proposition}[Envelope estimate]
\label{prop:envelope}
With $S(\tau) := \log\varphi_{\mathrm{sp}}(t) = g_\tau(w^*)+w^*-\tfrac12\log(2\pi V(r^*))$ and
$P(\tau) := Q(w_0)+B_0+w_0-\tfrac12\log\bigl(2\pi(B_0-\tfrac1c)\bigr)$,
\[
|S(\tau) - P(\tau) - H(w_0(\tau))| = O(1/\tau) \qquad (\tau\to\infty),
\]
uniformly in the phase of $\tau$ modulo $c$.
\end{proposition}

\begin{proof}[Sketch]
Expanding $S(\tau)-P(\tau)-H(w_0)$ exactly using $L=Q+H+\Delta$ and Lemma~\ref{lem:saddles}
produces four terms: the true envelope gap $g_{0,\tau}(w^*)-g_{0,\tau}(w_0)$, which is
$O(B_0\,(w^*-w_0)^2) = O(1/B_0)$; the periodic-part difference $H(w^*)-H(w_0) =
O(\eta\,|w^*-w_0|) = O(1/B_0)$ using \eqref{eq:Hbounds}; the saddle-location shift $w^*-w_0$
itself, $O(1/B_0)$; and a normalization term comparing $V(r^*)$ to $B_0 - 1/c$, again $O(1/B_0)$
using \eqref{eq:Hbounds}'s bound on $H''$. Since $\tau = w_0 - \log B_0$, $B_0/\tau \to 1/c$, and
all four terms are $O(1/\tau)$ with constants free of the phase of $\tau$.
\end{proof}

\section{The Berg--Kr\"uppel identity}
\label{sec:bk}

\begin{proposition}
\label{prop:bkidentity}
$P(\tau)$ is identical, as a function of $\tau$, to the exact saddlepoint expression Berg and
Kr\"uppel construct in the proof of their Proposition~9.1 -- the pre-asymptotic quantity their own
argument shows is asymptotic, as $t\to0^+$, to the explicit closed form $\varphi_0$ -- not merely
asymptotic to it. With $Y = cB_0$, $r = -ca-\log c$, the saddle equation for that construction is
exactly $Y-\log Y = \tau+r$, and $P(\tau) = Q(w_0)+B_0+w_0-\tfrac12\log(2\pi(B_0-1/c))$
coincides term for term with theirs: $2\beta y - \alpha = B_0$, $2\beta y - 2\beta - \alpha = B_0 -
1/c$, $\alpha y - \beta y^2 = Q(w_0)$ (in their notation, $y=w_0$). Consequently $P(\tau) -
\log\varphi_0(e^{-\tau}) \to 0$ (Section~\ref{sec:bk} below, following this proposition) is exactly
their own Proposition~9.1 applied to this identity, not a separate asymptotic claim about $P$.
\end{proposition}

\begin{proof}[Sketch]
Direct substitution, using $\alpha = a$, $\beta = 1/(2c)$ from their equation~(9.4) applied to
$a=3$, $\lambda=2/3$.
\end{proof}

Series-reverting $Y - \log Y = \tau + r$ to second order, $Y = \tau + \log\tau + r +
(\log\tau+r)/\tau + O(\tau^{-2}\log^2\tau)$, and substituting into $P$ gives
\[
P(\tau) = -\frac{(\tau+\log\tau)^2}{2c} + \Bigl(1+\tfrac1c-\tfrac{r}{c}\Bigr)\tau +
\Bigl(\tfrac12-\tfrac{r}{c}\Bigr)\log\tau + C_P + O\!\left(\frac{\log^2\tau}{\tau}\right),
\]
\[
C_P = -\frac{r^2}{2c}+r+ca+\frac{ca^2}{2}-\frac{\log2\pi}{2}+\frac{\log c}{2} =
-0.9576743183133982760669122497855855\ldots.
\]
Matching the $\tau$ and $\log\tau$ coefficients against Berg and Kr\"uppel's $\gamma$,
$\delta_{\mathrm{BK}}$ (their equations following~(9.6)) reproduces both exactly, and
\[
C_P = \varepsilon\log(2\beta) - \tfrac12\log(2\pi) = \log\frac{(2\beta)^\varepsilon}{\sqrt{2\pi}},
\]
so $P(\tau) - \log\varphi_{0,\mathrm{bare}}(e^{-\tau}) \to C_P$ and, using the prefactor-included
normalization, $P(\tau) - \log\varphi_0(e^{-\tau}) \to 0$ with rate $O(\tau^{-1}\log^2\tau)$: this
is precisely Berg and Kr\"uppel's own Proposition~9.1 (asymptotic, by construction, not exact)
applied to the identity of Proposition~\ref{prop:bkidentity}, not an independent asymptotic claim
about $P$ itself.

\begin{proof}[Proof of Corollary~\ref{cor:bk}]
By Proposition~\ref{prop:envelope} and \ref{prop:bkidentity}, $S(\tau) = \log\varphi_0(e^{-\tau}) +
H(w_0(\tau)) + O(\tau^{-1}\log^2\tau)$: the periodic factor multiplying Berg and Kr\"uppel's
asymptotic solution, in equation~(9.7) for this eigenfunction, is $H$, matching the product formula
their Proposition~9.3 supplies for it, as noted after Theorem~\ref{thm:H}. Proposition~
\ref{prop:Hfourier}'s Fourier series is the new, closed form; Proposition~\ref{prop:Hcert}'s
certificate is what neither their product formula nor their closing remark (that they only
\emph{expect} $\varphi$ bounded and bounded away from zero relative to $\varphi_0$) settles.
\end{proof}

\section{Proofs of the main theorems}
\label{sec:proofs}

\begin{proof}[Proof of Theorem~\ref{thm:conj3}]
Wirsching's class $\widetilde{A}_\delta$ (his equations~(1.5) and~(3.2)) consists of sequences
$(x_l)$ with $\lfloor 3^l x_l\rfloor = k_l$ satisfying $|l-k_l|\le\delta\sqrt l$ for every $l$; hence
$\lambda_l := 3^lx_l/l \to 1$ uniformly at rate $O(\delta/\sqrt l)$ on the class. Writing $z_l =
\lambda_l\, l\, 3^{-l}$, $\tau_l = -\log z_l = lc - \log l - \log\lambda_l$, the saddle location
satisfies $w_0(\tau_l) = lc - \log\lambda_l + O(1/l)$ by Lemma~\ref{lem:saddles}'s defining
equation, so $w_0(\tau_l) \bmod c \to 0$ uniformly on the class, at rate $O(1/\sqrt l)$. By
Theorem~\ref{thm:formulaA}, Proposition~\ref{prop:envelope}, and continuity of $H$ (indeed
Lipschitz, by \eqref{eq:Hbounds}),
\[
\log\varphi(z_l) - \log\varphi_0(z_l) = H(w_0(\tau_l)) + o(1) \to H(0)
\]
uniformly on $\widetilde{A}_\delta$, i.e. $\varphi(z_l)/\varphi_0(z_l) \to e^{H(0)}$ uniformly. This
is exactly Conjecture 3, with the constant given explicitly.
\end{proof}

\begin{proof}[Proof of Theorem~\ref{thm:unrestricted}]
As $t=e^{-\tau}$ ranges over $(0,1/2)$ with no restriction, $w_0(\tau) \bmod c$ ranges over a full
period (since $\Phi_0$ is a bijection and $\tau\mapsto\tau \bmod c$ composed with $\Phi_0^{-1}$ is
onto one period as $\tau\to\infty$), so by the same combination of results used above,
$\log\varphi(t)-\log\varphi_0(t) = H(w_0(\tau))+o(1)$ takes values arbitrarily close to both $\sup H$
and $\inf H$ infinitely often. By Proposition~\ref{prop:Hcert}, $\operatorname{osc}(H) > 0$
rigorously, so $\varphi(t)/\varphi_0(t)$ does not converge, and
$\limsup/\liminf \ge e^{\operatorname{osc}(H)} \ge 1.0004188$.
\end{proof}

\begin{proposition}[Wirsching's actual requirement]
\label{prop:wirschingreq}
Write $x_l^+$ for the upper endpoint of Wirsching's own covering interval at level $l$ (playing the
same role as this paper's $z_l$). Wirsching's own use of Conjecture 3, via
$\varphi'(x)=\tfrac92\varphi(3x)$ on $(0,\tfrac23)$ and his own calculation, his equation~(7.13),
that $\lim_l (2/3^{l+1})\varphi_0'(x_l)/\varphi_0(x_l) = 2/3$ uniformly on the class, reduces to
\[
\limsup_l \Lambda_l < \tfrac32, \qquad
\Lambda_l := \frac{\varphi(3x_l^+)/\varphi_0(3x_l^+)}{\varphi(x_l^+)/\varphi_0(x_l^+)}.
\]
This holds with a wide margin: $\Lambda_l \to 1$.
\end{proposition}

\begin{proof}
The phase shift under $t\mapsto3t$, i.e. $\tau\mapsto\tau-c$, is $w_0(\tau-c)-w_0(\tau)+c =
O(1/\tau)$ by Lemma~\ref{lem:saddles}'s defining equation, independently of the phase of $\tau$.
By Proposition~\ref{prop:envelope}, $\log\varphi(t)-\log\varphi_0(t) = H(w_0(\tau)) + O(1/\tau)$
uniformly in phase, so
\[
\log\Lambda_l = H(w_0(\tau_l-c)) - H(w_0(\tau_l)) + O(1/\tau_l) = O(1/\tau_l) \to 0,
\]
using that $H$ is Lipschitz (\eqref{eq:Hbounds}) and the phase shift above is $O(1/\tau_l)$; hence
$\Lambda_l \to 1$, well inside $\limsup \Lambda_l < 3/2$.
\end{proof}

\section{Discussion}
\label{sec:discussion}

These results settle Conjecture 3 alone. Wirsching's positive-predecessor-density argument reduces
to three conjectures; Conjectures 1 and 2, concerning respectively a pointwise-versus-average
transfer for the Elka functions and a uniformity-near-a-singularity statement for the limiting
operator $S_\infty$, are untouched here, and the full theorem is not established by this paper.

Wirsching stated the weaker Conjecture 3 instead of the stronger $\varphi\sim\kappa\varphi_0$ he
considers first. Theorem~\ref{thm:unrestricted} shows why: the stronger statement fails by a
mechanism with a certified, positive amplitude, not by some correctable technical looseness. The
class on which the weaker statement holds is exactly the one phase, among the continuum swept by
$H$, at which the ratio converges. That phase sits close to $H$'s minimum: $H(0) - \min H =
5.25\times10^{-6}$ against a total oscillation of $4.19\times10^{-4}$. We do not know whether this
is coincidental.

The rate is not established beyond $o(1)$. The sharper Edgeworth correction of
Remark~\ref{rem:edgeworth}, and a correspondingly sharper form of the envelope and Berg--Kr\"uppel
comparisons, would give an explicit polynomial rate throughout; neither theorem needs it.

\section{Code and data availability}
\label{sec:code}

Every certified or numerically-checked claim in this paper is backed by a runnable script in the
accompanying repository, \url{https://github.com/faculdade/wirsching-conjecture3-proof}, organized
by section with a README per script stating what it verifies and how to run it. This includes the
interval-arithmetic certificates behind Proposition~\ref{prop:Hcert} and the bounds
\eqref{eq:Hbounds}, the constant chain behind Theorem~\ref{thm:formulaA}, and the identity of
Proposition~\ref{prop:bkidentity}.

\subsection*{Acknowledgments}
AI assistance was used for textual review and translation.

\begin{thebibliography}{99}

\bibitem{wirsching2003}
G.~J. Wirsching, \emph{On the problem of positive predecessor density in $3N+1$ dynamics},
Discrete Contin. Dyn. Syst. \textbf{9} (2003), no.~3, 771--787.

\bibitem{bergkruppel1998}
L.~Berg and M.~Kr\"uppel, \emph{On the solution of an integral-functional equation with a
parameter}, Z. Anal. Anwendungen \textbf{17} (1998), no.~1, 159--181.

\bibitem{wirsching1998}
G.~J. Wirsching, \emph{The Dynamical System Generated by the $3n+1$ Function}, Lecture Notes in
Mathematics, vol.~1681, Springer, Berlin, 1998.

\bibitem{rvachev1990}
V.~A. Rvachev, \emph{Compactly supported solutions of functional-differential equations and their
applications}, Russian Math. Surveys \textbf{45} (1990), no.~1, 87--120.

\bibitem{debruijn1948}
N.~G. de Bruijn, \emph{On Mahler's partition problem}, Nederl. Akad. Wetensch., Proc.
\textbf{51} (1948), 659--669 (also Indag. Math. \textbf{10}, 210--220).

\bibitem{erdosrichmond1976}
P.~Erd\H{o}s and L.~B. Richmond, \emph{Concerning periodicity in the asymptotic behaviour of
partition functions}, J. Austral. Math. Soc. Ser.~A \textbf{21} (1976), 447--456.

\bibitem{katomcleod1971}
T.~Kato and J.~B. McLeod, \emph{The functional-differential equation $y'(x) = ay(\lambda x) +
by(x)$}, Bull. Amer. Math. Soc. \textbf{77} (1971), no.~6, 891--937.

\bibitem{derfel1995}
G.~Derfel, \emph{Functional-differential and functional equations with rescaling}, in Operator
Theory: Advances and Applications, vol.~80, Birkh\"auser, Basel, 1995, 100--111.

\end{thebibliography}

\end{document}
